\frontmatterchapter{Resumo}

Os sistemas ciberfísicos dependem cada vez mais de controlo autónomo distribuído, apesar de operarem com capacidades heterogéneas, configurações variáveis, informação incompleta, limitações de comunicação e restrições físicas. A Aprendizagem por Reforço Multiagente (MARL) oferece uma abordagem promissora para este contexto, mas um bom desempenho em simulação não demonstra, por si só, que uma política aprendida possa ser utilizada de forma fiável num sistema operacional. Esta tese aborda essa transição através de quatro requisitos relacionados: escalabilidade, capacidade de implantação, segurança operacional e robustez. As Comunidades de Energia Renovável e o carregamento de veículos elétricos constituem o principal domínio empírico, reunindo participantes distribuídos, ativos flexíveis heterogéneos, disponibilidade variável, objetivos locais e coletivos e infraestrutura física partilhada.

A investigação desenvolve o EnergAIze como uma arquitetura integrada para o desenvolvimento e operação de controlo inteligente de energia. Esta arquitetura liga dados e cenários, um simulador alargado derivado do CityLearn, desenvolvimento de controladores, experimentação reprodutível, governação de políticas, implantação distribuída, monitorização e recuperação. Neste contexto, uma progressão de métodos baseados em regras, otimização, aprendizagem por reforço e aprendizagem multiagente estabelece soluções progressivamente mais capazes e produz evidência sobre as condições em que cada uma se mantém eficaz.

Estes resultados motivam o TI-MARL, a principal contribuição algorítmica da tese. O TI-MARL considera problemas multiagente nos quais a interface de observação e ação pode mudar durante a operação. Um schema versionado de capacidades e informação de saúde em runtime são compilados numa interface tipada associada às entidades, contendo as observações, grupos de ação, portas de ação parametrizadas, dependências, condições de validade e limites atualmente disponíveis. Uma política descentralizada atua diretamente sobre esta interface variável, constrói uma representação local comum e produz um bundle coerente de ações para as entidades controláveis presentes. As restrições locais não negociáveis são incorporadas na construção das ações, enquanto a coordenação entre agentes é aprendida através de observações partilhadas autorizadas, objetivos acoplados, da transição física conjunta e de críticos centralizados capazes de representar populações variáveis durante o treino.

Os métodos e a infraestrutura são avaliados por verificação do simulador, comparação de controladores, reconfiguração dinâmica, perturbações e falhas, ablações, escalabilidade, compatibilidade, execução no edge e cenários energéticos integrados. Componentes selecionados são ainda avaliados com dados reais, plataformas externas, infraestrutura de carregamento e demonstradores físicos, distinguindo os níveis de evidência entre simulação e controlo operacional direto. No conjunto, a tese mostra como mecanismos algorítmicos, de representação, experimentação e ciclo de vida podem preservar controlo distribuído e fisicamente válido perante mudanças nos agentes, interfaces, condições e políticas implantadas.

\noindent\textbf{Palavras-chave:} Aprendizagem por Reforço Multiagente; Sistemas Ciberfísicos; Comunidades de Energia Renovável; TI-MARL; Escalabilidade; Segurança e Robustez; Ciclo de Vida de Políticas.
